\begin{center}\begin{tikzpicture}[font=\small,>=stealth,line join=round]
\draw[gray!40,rounded corners=3pt] (0,0) rectangle (14.2,11.6);
\node at (7.1,11.0) {例6：连续三次递减，最终大小关系相反};
\draw[->,gray] (1.0,9.6)--(7.5,9.6) node[right]{$x$};
\fill[blue] (2.3200000000000003,9.6) circle(2.2pt);
\node[above left=7pt,fill=white,inner sep=1pt] at (2.3200000000000003,9.6) {$2$};
\draw[blue,dashed,thick,->,shorten <=3pt,shorten >=3pt] (2.3200000000000003,9.6)--(5.98,7.0);
\fill[orange!85!black] (4.15,9.6) circle(2.2pt);
\node[above right=7pt,fill=white,inner sep=1pt] at (4.15,9.6) {$5$};
\draw[orange!85!black,dashed,thick,->,shorten <=3pt,shorten >=3pt] (4.15,9.6)--(4.15,7.0);
\draw[->,gray] (1.0,7.0)--(7.5,7.0) node[right]{$t$};
\fill[blue] (5.98,7.0) circle(2.2pt);
\node[above right=7pt,fill=white,inner sep=1pt] at (5.98,7.0) {$8$};
\draw[blue,dashed,thick,->,shorten <=3pt,shorten >=3pt] (5.98,7.0)--(2.3200000000000003,4.4);
\fill[orange!85!black] (4.15,7.0) circle(2.2pt);
\node[above left=7pt,fill=white,inner sep=1pt] at (4.15,7.0) {$5$};
\draw[orange!85!black,dashed,thick,->,shorten <=3pt,shorten >=3pt] (4.15,7.0)--(4.15,4.4);
\draw[->,gray] (1.0,4.4)--(7.5,4.4) node[right]{$u$};
\fill[blue] (2.3200000000000003,4.4) circle(2.2pt);
\node[above left=7pt,fill=white,inner sep=1pt] at (2.3200000000000003,4.4) {$2$};
\draw[blue,dashed,thick,->,shorten <=3pt,shorten >=3pt] (2.3200000000000003,4.4)--(5.98,1.8);
\fill[orange!85!black] (4.15,4.4) circle(2.2pt);
\node[above left=7pt,fill=white,inner sep=1pt] at (4.15,4.4) {$5$};
\draw[orange!85!black,dashed,thick,->,shorten <=3pt,shorten >=3pt] (4.15,4.4)--(4.15,1.8);
\draw[->,gray] (1.0,1.8)--(7.5,1.8) node[right]{$y$};
\fill[blue] (5.98,1.8) circle(2.2pt);
\node[below left=7pt,fill=white,inner sep=1pt] at (5.98,1.8) {$8$};
\fill[orange!85!black] (4.15,1.8) circle(2.2pt);
\node[below right=7pt,fill=white,inner sep=1pt] at (4.15,1.8) {$5$};
\node[align=left,anchor=west] at (8.4,6.0){每层都是 $10-$ 上一层\\[14pt]第一层：$2<5\to8>5$\\[12pt]第二层：$8>5\to2<5$\\[12pt]第三层：$2<5\to8>5$\\[14pt]三次反向，整体递减};
\node[text width=12.8cm,align=center] at (7.1,.65) {蓝点与橙点各自代表同一个输入；各行数轴向右增大，连接线追踪实际数值。};
\end{tikzpicture}\end{center}
